\documentclass[9pt,a4paper,twocolumn]{ctexart}

% --- Packages ---
\usepackage[top=1.2cm, bottom=1.2cm, left=1.2cm, right=1.2cm, columnsep=0.8cm]{geometry}
\usepackage{hyperref}
\usepackage{graphicx}
\usepackage{float}
\usepackage{longtable}
\usepackage{tabularx}
\usepackage{booktabs}
\usepackage{enumitem}
\usepackage{xcolor}
\usepackage{fancyhdr}
\usepackage{listings}
\usepackage{fontspec}
\usepackage{titlesec}
\usepackage{caption}

% --- Code block styling ---
\definecolor{codebg}{RGB}{245,245,245}
\definecolor{codeframe}{RGB}{220,220,220}
\definecolor{codekw}{RGB}{0,0,192}
\definecolor{codecomment}{RGB}{0,128,0}
\definecolor{codestring}{RGB}{163,21,21}
\definecolor{codenum}{RGB}{128,0,128}
\definecolor{codepunct}{RGB}{90,90,90}
\definecolor{badgegreen}{RGB}{34,139,34}
\definecolor{badgeblue}{RGB}{30,144,255}
\definecolor{badgeblack}{RGB}{50,50,50}

% --- Difficulty badges (rounded corners via tikz, pre-rendered in saveboxes) ---
\usepackage{tikz}
\newsavebox{\badgechubox}
\savebox{\badgechubox}{\tikz[baseline]{\node[fill=badgegreen,text=white,rounded corners=2.5pt,inner xsep=4pt,inner ysep=1.5pt,font=\sffamily\footnotesize\bfseries]{初级};}}
\newsavebox{\badgezhongbox}
\savebox{\badgezhongbox}{\tikz[baseline]{\node[fill=badgeblue,text=white,rounded corners=2.5pt,inner xsep=4pt,inner ysep=1.5pt,font=\sffamily\footnotesize\bfseries]{中级};}}
\newsavebox{\badgegaobox}
\savebox{\badgegaobox}{\tikz[baseline]{\node[fill=badgeblack,text=white,rounded corners=2.5pt,inner xsep=4pt,inner ysep=1.5pt,font=\sffamily\footnotesize\bfseries]{高级};}}
\newcommand{\lvchu}{\usebox{\badgechubox}}
\newcommand{\lvzhong}{\usebox{\badgezhongbox}}
\newcommand{\lvgao}{\usebox{\badgegaobox}}
\pdfstringdefDisableCommands{%
  \def\lvchu{[初级]}%
  \def\lvzhong{[中级]}%
  \def\lvgao{[高级]}%
}

\lstset{
  basicstyle=\ttfamily\scriptsize,
  keywordstyle=\color{codekw}\bfseries,
  commentstyle=\color{codecomment}\itshape,
  stringstyle=\color{codestring},
  backgroundcolor=\color{codebg},
  frame=single,
  rulecolor=\color{codeframe},
  breaklines=true,
  breakatwhitespace=false,
  breakindent=0pt,
  postbreak=\mbox{\textcolor{red}{$\hookrightarrow$}\space},
  numbers=left,
  numberstyle=\tiny\color{gray},
  columns=flexible,
  keepspaces=true,
  showstringspaces=false,
  tabsize=2,
  aboveskip=4pt,
  belowskip=4pt,
  extendedchars=true,
  literate={，}{{，}}1 {；}{{；}}1 {！}{{！}}1 {？}{{？}}1 {“}{{“}}1 {”}{{”}}1 {（}{{（}}1 {）}{{）}}1 {《}{{《}}1 {》}{{》}}1,
}

% --- Extra language definitions (no built-in listings support) ---
\lstdefinelanguage{json}{
  morekeywords={true,false,null},
  sensitive=true,
  morestring=[b]",
  comment=[l]{//},
  morecomment=[s]{/*}{*/},
}
\lstdefinelanguage{yaml}{
  morekeywords={true,false,null,yes,no,on,off},
  sensitive=false,
  morecomment=[l]{\#},
  morestring=[b]",
  morestring=[d]',
}
\lstdefinelanguage{http}{
  morekeywords={GET,POST,PUT,DELETE,PATCH,HEAD,OPTIONS,HTTP,Host,Accept,Authorization,Content-Type,Content-Length,User-Agent},
  sensitive=true,
  morecomment=[l]{\#},
  morestring=[b]",
}
\lstdefinelanguage{TypeScript}{
  morekeywords={abstract,any,as,async,await,break,case,catch,class,const,continue,debugger,declare,default,delete,do,else,enum,export,extends,false,finally,for,from,function,get,if,implements,import,in,instanceof,interface,keyof,let,module,namespace,never,new,null,number,of,package,private,protected,public,readonly,return,set,static,string,super,switch,symbol,this,throw,true,try,type,typeof,undefined,unknown,var,void,while,with,yield,Record,Array,Promise,AsyncIterable,ReadableStream,Date,Error,Object,JSON},
  sensitive=true,
  morecomment=[l]{//},
  morecomment=[s]{/*}{*/},
  morestring=[b]",
  morestring=[b]',
  morestring=[b]`{`},
}

% --- Compact spacing ---
\linespread{0.95}
\setlength{\parskip}{1pt plus 1pt minus 1pt}
\setlength{\parindent}{0pt}

% --- Table styling ---
\renewcommand{\arraystretch}{1.1}

% --- Heading styling (numbered) ---
\titleformat{\section}{\normalsize\bfseries}{\thesection\quad}{0em}{}
\titlespacing{\section}{0pt}{6pt}{2pt}
\titleformat{\subsection}{\small\bfseries}{\thesubsection\quad}{0em}{}
\titlespacing{\subsection}{0pt}{4pt}{1pt}
\titleformat{\subsubsection}{\small\bfseries}{\thesubsubsection\quad}{0em}{}
\titlespacing{\subsubsection}{0pt}{3pt}{1pt}

% --- Page style ---
\pagestyle{fancy}
\fancyhf{}
\fancyhead[L]{\small\emph{\leftmark}}
\fancyhead[R]{\small\thepage}
\renewcommand{\headrulewidth}{0.4pt}
\renewcommand{\sectionmark}[1]{}  % prevent sections from overwriting leftmark

% --- Hyperref setup ---
\hypersetup{
  colorlinks=true,
  linkcolor=blue,
  urlcolor=blue,
  citecolor=blue,
}

% --- Caption format ---
\DeclareCaptionFormat{myformat}{\small\bfseries #1#2#3}
\captionsetup{format=myformat}

\begin{document}

\title{AI Agent 面试题总结}
\date{}
\maketitle
\markboth{TypeScript 版 Agent 知识}{ TypeScript 版 Agent 知识}

\section*{面试官真正想考什么}


Agent 题本质上在考“复杂 AI 应用怎么编排”。可以按下面几个层次准备。


{\footnotesize
\begin{tabular}{|p{\dimexpr 0.216\columnwidth-2\tabcolsep\relax}|p{\dimexpr 0.398\columnwidth-2\tabcolsep\relax}|p{\dimexpr 0.386\columnwidth-2\tabcolsep\relax}|}
\hline
\textbf{考察方向} & \textbf{面试官想确认什么} & \textbf{常见扣分点} \\
\hline
Agent 基础 & 你能否讲清 Agent、Workflow、普通 Chatbot 的区别 & 把 Agent 说成“会自动思考的机器人” \\
\hline
Agent Loop & 你是否理解推理、行动、观察、修正的循环 & 只讲工具调用，不讲观察和迭代 \\
\hline
Context Engineering & 你是否知道上下文质量决定 Agent 表现 & 只会调 Prompt，不会管理上下文 \\
\hline
Memory & 你是否能区分短期状态、长期事实和经验沉淀 & 把历史聊天记录等同于记忆系统 \\
\hline
Tools/MCP/Skills & 你是否知道工具接入、调用意图和任务 SOP 的边界 & 把 MCP、Function Calling、Skills 混为一谈 \\
\hline
Workflow/Harness & 你是否具备生产级 Agent 工程化思维 & 盲目追求纯 Agent，不考虑可控性 \\
\hline
\end{tabular}
}



回答 Agent 题时，建议少讲“智能”，多讲“约束”。因为真实项目里，Agent 最大的问题不是不会做事，而是不稳定、不可控、难排查、成本高。



\section*{Agent 基础}


参考文章：《AI Agent 核心概念：Agent Loop、Context Engineering、Tools 注册》



这一组题是 Agent 面试的入口。重点不是背公式，而是讲清 Agent 和传统程序、Workflow 的边界。



建议掌握这些关键点：


\begin{itemize}[itemsep=1pt, topsep=2pt]
  \item Agent 可以理解为 LLM + Planning + Memory + Tools 的组合，但这个公式只是起点，不是完整生产架构。
  \item 普通 Chatbot 主要回答问题，Agent 更强调多步骤任务执行和外部工具调用。
  \item Workflow 的路径更固定，适合流程清晰、需要可控性的场景；纯 Agent 更适合路径难提前穷举的开放任务。
  \item ReAct、Plan-and-Execute、Reflection、Multi-Agent 不是越复杂越好，要结合任务复杂度、调试成本和容错要求选择。
\end{itemize}


高频面试题：


\begin{itemize}[itemsep=1pt, topsep=2pt]
  \item AI Agent 是什么？和普通 Chatbot 有什么区别？
  \item Agent = LLM + Planning + Memory + Tools 这条公式怎么理解？
  \item Agent Loop 的完整流程是什么？
  \item Agent 和传统编程、Workflow 的核心区别是什么？
  \item ReAct、Plan-and-Execute、Reflection、Multi-Agent 分别适合什么场景？
  \item Tools 注册时，工具 description 为什么很关键？
  \item 什么时候用纯 Agent，什么时候用 Workflow 或 Agentic Workflow？
  \item Multi-Agent 协作的主要问题是什么？为什么生产里不能盲目上多 Agent？
\end{itemize}


一个更稳的回答方式是：先承认 Agent 的动态决策能力，再补上它的代价。比如纯 Agent 灵活，但调试难、轨迹不稳定、Token 成本高；Workflow 可控，但前期流程拆解要求高。To B 场景通常会优先选择 Workflow 或 Agentic Workflow，把关键路径控制住，只在必要节点让模型做判断。



\begin{figure}[H]
\centering
\includegraphics[width=0.85\textwidth]{/home/zhang/demos/latex_demos/Agent知识/HTML/AI Agent 面试题总结 _ JavaGuide_files/agent-core-arch.jpg}
\caption*{AI Agent 核心架构}
\end{figure}



\begin{figure}[H]
\centering
\includegraphics[width=0.85\textwidth]{/home/zhang/demos/latex_demos/Agent知识/HTML/AI Agent 面试题总结 _ JavaGuide_files/agent-loop-flow.jpg}
\caption*{Agent Loop 工作流程}
\end{figure}



\section*{Agent Memory}


参考文章：《AI Agent 记忆系统：短期记忆、长期记忆与记忆演化机制》



Memory 题经常被问得很细，因为它能区分“玩过 Demo”和“做过系统”的候选人。真正的记忆系统不是把聊天记录一股脑塞回上下文，而是对信息进行分层、筛选、压缩、更新和治理。



建议掌握这些关键点：


\begin{itemize}[itemsep=1pt, topsep=2pt]
  \item 短期记忆更像当前任务状态，负责记录这一轮任务里必须保留的信息。
  \item 长期记忆更像跨会话知识，负责沉淀用户偏好、团队规则、历史决策和经验。
  \item 向量记忆适合语义检索，Markdown 记忆适合规则、偏好、项目约定这类可读可审查的信息。
  \item 记忆写入不能完全放任模型自动决定，否则容易写入错误、过时、重复或敏感信息。
  \item 团队共享记忆最好走 Git、PR 和 Review，便于审计和回滚。
\end{itemize}


高频面试题：


\begin{itemize}[itemsep=1pt, topsep=2pt]
  \item Agent 的短期记忆和长期记忆有什么区别？
  \item Agent 记忆系统要解决哪些核心问题？
  \item 向量记忆和 Markdown 记忆分别适合什么场景？
  \item Auto Memory 是什么？它为什么不能无限自动写入？
  \item 团队共享记忆为什么适合走 Git 和 Code Review？
  \item 记忆压缩、记忆过期、记忆冲突应该怎么处理？
  \item 如何避免长期记忆污染上下文？
  \item 面试里怎么讲“有记忆”不是简单保存聊天记录？
\end{itemize}


如果被追问“怎么设计记忆系统”，可以按读写链路回答：先定义哪些信息允许写入，再做敏感信息过滤和去重；写入时记录来源、时间、置信度和作用域；读取时根据任务检索相关记忆，而不是全量注入；过期或冲突时通过人工审核或规则策略处理。



\begin{figure}[H]
\centering
\includegraphics[width=0.85\textwidth]{/home/zhang/demos/latex_demos/Agent知识/HTML/AI Agent 面试题总结 _ JavaGuide_files/agent-memory-memory-taxonomy.pdf}
\caption*{Agent 记忆分类全景图}
\end{figure}



\section*{Prompt 与 Context Engineering}


参考文章：《大模型提示词工程实践指南》、《上下文工程实战指南：让 Agent 少犯蠢的工程方法论》



Agent 场景下，Prompt 只是入口，Context 才是持续影响模型行为的“工作台”。很多 Agent 不稳定，不是 Prompt 写得不够长，而是上下文里噪声太多、关键约束位置太差、工具结果格式混乱、历史状态没有结构化。



建议掌握这些关键点：


\begin{itemize}[itemsep=1pt, topsep=2pt]
  \item Prompt Engineering 关注指令怎么写清楚，Context Engineering 关注什么信息在什么时机进入模型窗口。
  \item Agent 上下文通常包含系统规则、任务目标、历史状态、工具说明、工具结果、用户偏好、检索证据和中间计划。
  \item 长任务要做上下文压缩、结构化笔记、任务状态持久化和必要的 Sub-agent 拆分。
  \item Prompt 注入不能只靠提醒模型“不要听用户恶意指令”，还要靠权限隔离、工具白名单、输出校验和审计。
\end{itemize}


高频面试题：


\begin{itemize}[itemsep=1pt, topsep=2pt]
  \item Prompt Engineering 和 Context Engineering 有什么区别？
  \item Prompt 四要素 Role、Task、Context、Format 分别解决什么问题？
  \item Few-Shot、CoT、任务分解、结构化输出分别适合什么场景？
  \item Prompt 注入攻击是什么？常见防护方式有哪些？
  \item 为什么 Agent 场景下只优化 Prompt 不够？
  \item Context Engineering 要解决哪些问题？
  \item 静态规则、动态信息、工具结果、记忆应该如何进入上下文？
  \item 长任务上下文溢出时，Compaction、结构化笔记、Sub-agent 分别怎么用？
\end{itemize}


答这类题时，可以抓住一句话：\textbf{Prompt 决定模型收到什么指令，Context 决定模型实际看到什么世界。} Agent 一旦进入多轮工具调用，后者往往更重要。



\begin{figure}[H]
\centering
\includegraphics[width=0.85\textwidth]{/home/zhang/demos/latex_demos/Agent知识/HTML/AI Agent 面试题总结 _ JavaGuide_files/context-engineering-vs-prompt-engineering.jpg}
\caption*{Prompt engineering vs. context engineering}
\end{figure}



\section*{MCP 与 Agent Skills}


参考文章：《深入理解 MCP 协议：一次开发，多处复用》、《Agent Skills 是什么？和 Prompt、MCP 到底差在哪？》



这一组题考的是工具生态和能力复用。很多人会把 MCP、Function Calling、Skills 都说成“工具调用”，这样答会显得边界不清。



建议掌握这些关键点：


\begin{itemize}[itemsep=1pt, topsep=2pt]
  \item Function Calling 解决的是模型如何输出结构化工具调用意图。
  \item MCP 解决的是工具如何被标准化发现、描述、调用和返回结果。
  \item Skills 解决的是 Agent 做某类任务时，应该按什么经验和流程执行。
  \item MCP 更像能力接口，Skills 更像任务 SOP。二者可以组合使用。
  \item 生产级工具接入必须有权限、参数校验、审计、超时、重试和降级策略。
\end{itemize}


高频面试题：


\begin{itemize}[itemsep=1pt, topsep=2pt]
  \item MCP 解决什么问题？为什么常被类比成 AI 领域的 USB-C？
  \item MCP Client、MCP Server、Host 分别是什么？
  \item MCP 的 Tools、Resources、Prompts 分别解决什么问题？
  \item MCP 和 Function Calling 有什么区别？
  \item 生产级 MCP Server 要做哪些安全治理？
  \item Agent Skills 是什么？它和 Prompt、MCP、Function Calling 的边界是什么？
  \item Skills 为什么要延迟加载？
  \item Skill 路由怎么做？为什么它和 RAG 相似但目标不同？
  \item 写一个 \texttt{SKILL.md} 最容易踩哪些坑？
\end{itemize}


面试里可以这样概括：Function Calling 是“模型怎么表达要调工具”，MCP 是“工具怎么接入宿主”，Skills 是“Agent 做这类任务时按什么经验执行”。三者不是替代关系，而是不同层次的组合。



\section*{Harness Engineering}


参考文章：《一文搞懂 Harness Engineering：六层架构、上下文管理与一线团队实战》



Harness Engineering 是 Agent 面试里比较进阶的一块。它的核心思想是：不要把 Agent 表现完全归因于模型本身，模型之外的任务管理、上下文供给、工具反馈、验证机制、错误恢复，同样决定系统上限。



建议掌握这些关键点：


\begin{itemize}[itemsep=1pt, topsep=2pt]
  \item Agent = Model + Harness。模型负责推理和生成，Harness 负责把任务、上下文、工具和反馈组织起来。
  \item Harness 里的每个组件，本质上都编码了一个假设：模型单独做不好什么。
  \item 模型能力升级后，Harness 也要重新评估。有些过去必要的补丁，可能会变成新的复杂度。
  \item 上下文污染、代码熵积累、工具调用可靠性，是一线 Agent 工程里很常见的三类问题。
\end{itemize}


高频面试题：


\begin{itemize}[itemsep=1pt, topsep=2pt]
  \item Harness Engineering 是什么？它和 Prompt Engineering、Context Engineering 有什么关系？
  \item 为什么说 Agent = Model + Harness？
  \item Harness 的六层架构分别解决什么问题？
  \item 模型能力升级后，Harness 里的某些机制为什么需要重新验证？
  \item 上下文污染、代码熵积累、工具调用可靠性分别怎么治理？
  \item Agent 工程里为什么需要评测器、验证器和任务状态管理？
  \item 一线团队做 Agent 工程化时，共同遇到的难点是什么？
\end{itemize}


回答时别把 Harness 讲成新名词堆砌。更好的方式是用具体问题带出来：Agent 长任务中途跑偏，需要任务状态和阶段性检查；工具返回错误，模型需要可修复的错误反馈；代码生成重复实现已有逻辑，需要检索和去重机制。这些都是 Harness 要补的系统能力。



\begin{figure}[H]
\centering
\includegraphics[width=0.85\textwidth]{/home/zhang/demos/latex_demos/Agent知识/HTML/AI Agent 面试题总结 _ JavaGuide_files/harness-engineering-layers-arch.jpg}
\caption*{Harness 和 Prompt/Context Engineering 的关系}
\end{figure}



\section*{Workflow、Graph 与 Loop}


参考文章：《AI 工作流中的 Workflow、Graph 与 Loop：从概念到实现》



这一组题适合用来展示工程判断。很多业务场景并不适合纯 Agent，而是更适合把流程设计成 Graph，让模型只在必要节点做生成、判断或路由。



建议掌握这些关键点：


\begin{itemize}[itemsep=1pt, topsep=2pt]
  \item Workflow 是任务过程，Graph 是结构载体，Loop 是控制模式。
  \item Graph 中 Node 负责执行，Edge 负责流转，State 负责保存跨节点上下文。
  \item Loop 必须有继续条件、退出条件和安全边界，否则很容易死循环或烧 Token。
  \item State 更新要设计策略：单值字段 Replace，日志类字段 Append，并行写入字段需要 Reducer。
\end{itemize}


高频面试题：


\begin{itemize}[itemsep=1pt, topsep=2pt]
  \item 为什么 AI 系统需要工作流？
  \item Workflow、Graph、Loop 三者是什么关系？
  \item Graph Loop 和 Agent Loop 有什么区别？
  \item Loop 如何防止死循环？
  \item State 的更新策略怎么选？Replace、Append、Reducer 分别适合什么字段？
  \item 条件边和动态路由有什么区别？
  \item 工作流中断后怎么恢复？
  \item 工作流有哪些特有的安全风险？
\end{itemize}


面试官如果问“你会怎么设计一个复杂 Agent 流程”，可以先画出固定主链路，再说明哪些节点由模型判断，哪些节点必须由规则和代码控制。这样比直接说“让 Agent 自己规划”可信得多。



\section*{答题框架}


Agent 题可以用这条主线来回答：


\begin{enumerate}[itemsep=1pt, topsep=2pt]
  \item 先定义任务类型：是问答、检索、工具调用、多步骤任务，还是长周期任务。
  \item 再选择编排方式：纯 Agent、Workflow、Agentic Workflow 或 Multi-Agent。
  \item 接着讲核心组件：Context、Memory、Tools、MCP、Skills、State。
  \item 然后讲安全和稳定性：权限、校验、超时、重试、审计、成本控制。
  \item 最后讲评测：任务完成率、工具调用准确率、轨迹质量和失败样本回放。
\end{enumerate}


这个框架的好处是，它能把“Agent 很智能”拉回到“系统怎么设计”。



\section*{常见扣分点}

\begin{itemize}[itemsep=1pt, topsep=2pt]
  \item 把 Agent 讲成万能自动化，忽略失败恢复和安全边界。
  \item 只讲 Prompt，不讲上下文供给、工具结果和状态管理。
  \item 把 Memory 等同于历史聊天记录。
  \item 把 MCP、Function Calling、Skills 混成一个概念。
  \item 盲目推 Multi-Agent，不考虑通信成本、调试成本和一致性问题。
  \item 不知道什么时候该用 Workflow，而不是纯 Agent。
\end{itemize}


\section*{复习建议}


建议按这个顺序复习：


\begin{enumerate}[itemsep=1pt, topsep=2pt]
  \item 先看 Agent 基础，讲清 Agent、Chatbot、Workflow 的区别。
  \item 再看 Memory 和 Context Engineering，理解 Agent 稳定性的关键。
  \item 接着看 MCP、Skills、Function Calling，掌握工具生态边界。
  \item 最后看 Harness Engineering 和 Workflow，把知识收敛到生产级架构。
\end{enumerate}


复习时不要只问“Agent 是什么”，要继续追问：它如何拿到信息？如何调用工具？如何记住状态？如何失败恢复？如何评测？这些问题答清楚，才像真的做过 Agent。


\end{document}
